\documentclass[12pt]{article}
\usepackage[utf8]{inputenc}
\usepackage[T1]{fontenc}
\usepackage[margin=1in]{geometry}
\usepackage{setspace}
\usepackage{xcolor}
\usepackage[colorlinks=true, urlcolor=blue!60!black, linkcolor=blue!60!black]{hyperref}
\setlength{\parindent}{2em}
\setlength{\parskip}{1.25em}
\renewcommand{\baselinestretch}{1.0}
\begin{document}
\begin{center}
{\LARGE \textit{Washington v. Davis}, 426 U.S. 229 (1976)\par}
\vspace{0.8em}
{\large Case Brief\par}
\end{center}
\vspace{1em}

\section{Memory Jogger}

A facially neutral government policy that produces a racially disproportionate impact does not violate the equal protection component of the Fifth Amendment or the Equal Protection Clause of the Fourteenth Amendment unless the challenger proves a racially discriminatory purpose (\href{https://www.courtlistener.com/opinion/109469/washington-v-davis/}{CourtListener}; \href{https://supreme.justia.com/cases/federal/us/426/229/}{Justia}).

\section{Detailed Case Facts}

On April 10, 1970, two Black police officers, including Alfred E. Davis, filed suit in federal court against District of Columbia officials---including then-Commissioner (later Mayor) Walter E. Washington, the Chief of the Metropolitan Police Department, and the Commissioners of the United States Civil Service Commission---alleging racially discriminatory promotion policies within the District of Columbia Metropolitan Police Department (\href{https://supreme.justia.com/cases/federal/us/426/229/}{Justia}; \href{https://law.justia.com/cases/federal/appellate-courts/F2/512/956/286050/}{D.C. Cir. Opinion}). Two rejected Black applicants for police officer positions, George Harley and John D. Sellers, were permitted to intervene, challenging the Department's recruiting procedures as racially discriminatory under the Due Process Clause of the Fifth Amendment (which applies to the federal District of Columbia and incorporates equal protection principles under \textit{Bolling v. Sharpe}, 347 U.S. 497 (1954)), 42 U.S.C. \S~1981, and D.C. Code \S~1-320 (\href{https://supreme.justia.com/cases/federal/us/426/229/}{Justia}; \href{https://www.law.cornell.edu/constitution-conan/amendment-14/facially-neutral-laws-implicating-suspect-classifications}{Cornell LII}).

To be accepted into the Department's 17-week recruit training program, an applicant had to meet physical and character standards, possess a high school diploma or equivalent, and score at least 40 out of 80 on ``Test 21,'' a written personnel examination developed by the U.S. Civil Service Commission for general use throughout the federal civil service (\href{https://supreme.justia.com/cases/federal/us/426/229/}{Justia}). Test 21 measured verbal ability, vocabulary, reading, and comprehension (\href{https://supreme.justia.com/cases/federal/us/426/229/}{Justia}).

The intervenors did not allege intentional discrimination or purposeful discriminatory acts in the design or administration of Test 21; instead, they argued that the examination bore no demonstrated relationship to actual on-the-job police performance and excluded a disproportionately high percentage of Black applicants---four times as many Black applicants failed Test 21 as white applicants (\href{https://supreme.justia.com/cases/federal/us/426/229/}{Justia}; \href{https://www.oyez.org/cases/1975/74-1492}{Oyez}). At the same time, the undisputed record showed that since August 1969, 44\% of new police recruits had been Black---roughly matching both the proportion of Black officers on the force and the share of 20- to 29-year-old Black residents within the Department's 50-mile recruiting radius---and that the Department had systematically and affirmatively recruited Black candidates (\href{https://supreme.justia.com/cases/federal/us/426/229/}{Justia}). Furthermore, a Civil Service Commission validation study by D. L. Futransky and expert affidavits showed a significant positive correlation between Test 21 scores and final examination averages in the police recruit training program (\href{https://supreme.justia.com/cases/federal/us/426/229/}{Justia}).

\section{Procedural History}

\begin{itemize}
\item \textbf{United States District Court for the District of Columbia (\textit{Davis v. Washington}, 348 F. Supp. 15 (D.D.C. 1972)):} Focusing solely on Test 21 on cross-motions for summary judgment (with the respondents' motion resting exclusively on the Fifth Amendment Due Process Clause and the defendants' motions asserting compliance with both constitutional and statutory standards), the District Court found that Black applicants failed Test 21 at a higher rate and that the test had not been validated for subsequent on-the-job performance (\href{https://supreme.justia.com/cases/federal/us/426/229/}{Justia}). However, the court granted summary judgment to the District of Columbia and federal defendants, concluding that the Department's affirmative recruitment efforts, the 44\% Black recruit share, and Test 21's direct relationship to training-academy performance negated any inference of racial discrimination (\href{https://supreme.justia.com/cases/federal/us/426/229/}{Justia}).
\item \textbf{United States Court of Appeals for the District of Columbia Circuit (\textit{Davis v. Washington}, 168 U.S. App. D.C. 42, 512 F.2d 956 (1975)):} Over a dissent, the D.C. Circuit reversed and directed partial summary judgment for respondents (\href{https://law.justia.com/cases/federal/appellate-courts/F2/512/956/286050/}{D.C. Cir. Opinion}; \href{https://supreme.justia.com/cases/federal/us/426/229/}{Justia}). Importing the Title VII disparate-impact framework of \textit{Griggs v. Duke Power Co.}, 401 U.S. 424 (1971), into the Fifth Amendment constitutional inquiry, the Court of Appeals held that the absence of discriminatory intent was irrelevant, that the four-to-one failure disparity established a constitutional violation unless the government proved Test 21 predicted actual job performance, and that correlation with training-school exams was insufficient (\href{https://supreme.justia.com/cases/federal/us/426/229/}{Justia}).
\item \textbf{United States Supreme Court (\textit{Washington v. Davis}, 426 U.S. 229 (1976)):} Certiorari was granted at 423 U.S. 820 (1975) on the District of Columbia officials' petition (\href{https://supreme.justia.com/cases/federal/us/426/229/}{Justia}). Argued March 1, 1976, and decided June 7, 1976, the Supreme Court reversed the D.C. Circuit (\href{https://www.oyez.org/cases/1975/74-1492}{Oyez}; \href{https://www.courtlistener.com/opinion/109469/washington-v-davis/}{CourtListener}).
\end{itemize}

\section{Judicial Votes}

\textbf{7--2 decision} reversing the Court of Appeals (\href{https://www.oyez.org/cases/1975/74-1492}{Oyez}; \href{https://supreme.justia.com/cases/federal/us/426/229/}{Justia}):
\begin{itemize}
\item \textbf{Majority Opinion (Parts I, II, and III):} Justice Byron R. White delivered the opinion of the Court, joined in full by Chief Justice Warren E. Burger and Justices Harry A. Blackmun, Lewis F. Powell, Jr., William H. Rehnquist, and John Paul Stevens (with Justice Stevens noting a reservation as to the merits of the lower-court cases listed in footnote 12) (\href{https://supreme.justia.com/cases/federal/us/426/229/}{Justia}).
\item \textbf{Partial Joinder:} Justice Potter Stewart joined Parts I and II of the Court's opinion (the constitutional holding), declining to join Part III (the statutory validation holding) (\href{https://supreme.justia.com/cases/federal/us/426/229/}{Justia}).
\item \textbf{Concurring Opinion:} Justice John Paul Stevens filed a concurring opinion addressing both the constitutional intent standard and the statutory job-relatedness issue (\href{https://supreme.justia.com/cases/federal/us/426/229/}{Justia}).
\item \textbf{Dissenting Opinion:} Justice William J. Brennan, Jr., filed a dissenting opinion, joined by Justice Thurgood Marshall, contesting the Court's statutory holding in Part III without reaching the constitutional question in Part II (\href{https://supreme.justia.com/cases/federal/us/426/229/}{Justia}).
\end{itemize}

\section{Holding}

\textbf{Judgment:} Reversed (\href{https://supreme.justia.com/cases/federal/us/426/229/}{Justia}). The District Court properly granted summary judgment to the District of Columbia and federal defendants on both constitutional and statutory grounds (\href{https://supreme.justia.com/cases/federal/us/426/229/}{Justia}).

\textbf{Rule Announced:} First, under the equal protection component of the Fifth Amendment's Due Process Clause and the Fourteenth Amendment's Equal Protection Clause, a facially neutral law or official act is not unconstitutional solely because it produces a racially disproportionate impact; a challenger must establish a racially discriminatory purpose to trigger strict scrutiny, though disproportionate impact remains relevant circumstantial evidence from which discriminatory purpose may be inferred (\href{https://supreme.justia.com/cases/federal/us/426/229/}{Justia}; \href{https://www.law.cornell.edu/constitution-conan/amendment-14/facially-neutral-laws-implicating-suspect-classifications}{Cornell LII}). Absent discriminatory purpose, a facially neutral measure is reviewed under rational basis scrutiny (\href{https://supreme.justia.com/cases/federal/us/426/229/}{Justia}). Second, under the applicable federal civil service and District of Columbia statutory standards (which the parties treated as analogous to Title VII), an employment test may be validated by demonstrating a positive relationship to performance in a job-relevant training program, without necessarily proving direct correlation to post-training job performance (\href{https://supreme.justia.com/cases/federal/us/426/229/}{Justia}).

\textbf{Holding vs. Dicta:} The twin holdings are (1) that the Court of Appeals committed plain error under Supreme Court Rule 40(1)(d)(2) by importing Title VII's disparate-impact standard into Fifth Amendment equal protection analysis, and that Test 21 was constitutional because it was facially neutral, served a legitimate government interest in upgrading verbal competence, and lacked discriminatory purpose; and (2) that Test 21 satisfied the applicable non-Title VII statutory standards through training-program validation (\href{https://supreme.justia.com/cases/federal/us/426/229/}{Justia}). By contrast, Justice White's sweeping warning that a constitutional disparate-impact rule would jeopardize a broad array of tax, welfare, public service, regulatory, and licensing statutes---and his footnote 12 disapproval of numerous lower-court decisions involving urban renewal, zoning, public housing, and municipal services, including \textit{Metropolitan Housing Development Corp. v. Village of Arlington Heights}, then pending before the Court---went beyond the employment test at issue and functioned as influential dicta (as Justice Stevens's reservation and Justice Brennan's footnote 1 critique underscored) (\href{https://supreme.justia.com/cases/federal/us/426/229/}{Justia}).

\section{Analysis of Opinions}

\subsection{Majority Opinion (Justice White)}

In Part II, Justice White invoked Supreme Court Rule 40(1)(d)(2) to notice ``plain error'' because the D.C. Circuit had collapsed the Fifth Amendment standard into Title VII's statutory disparate-impact rule from \textit{Griggs v. Duke Power Co.} (\href{https://supreme.justia.com/cases/federal/us/426/229/}{Justia}). Surveying a century of equal protection precedent---including jury-exclusion cases (\textit{Strauder v. West Virginia}, \textit{Akins v. Texas}, \textit{Alexander v. Louisiana}), legislative apportionment (\textit{Wright v. Rockefeller}), school desegregation (\textit{Keyes v. School District No. 1}, distinguishing intentional \textit{de jure} segregation from \textit{de facto} segregation), and welfare benefits (\textit{Jefferson v. Hackney})---Justice White concluded that the Court's doctrine had consistently required purposeful discrimination (\href{https://supreme.justia.com/cases/federal/us/426/229/}{Justia}).

Justice White clarified that discriminatory purpose need not appear on the face of a statute and may be inferred from the ``totality of the relevant facts,'' including severe disproportionate impact (as in \textit{Yick Wo v. Hopkins} or systematic jury exclusion) (\href{https://supreme.justia.com/cases/federal/us/426/229/}{Justia}). He distinguished \textit{Palmer v. Thompson} (closing municipal swimming pools) and \textit{Wright v. Council of City of Emporia} (splitting a school district under a desegregation decree) as narrow decisions that did not displace the purpose requirement (\href{https://supreme.justia.com/cases/federal/us/426/229/}{Justia}). Applying the rule to Test 21, Justice White reasoned that the government may legitimately seek to upgrade employees' verbal and communicative skills, that Black applicants who failed Test 21 stood on the same constitutional footing as white applicants who failed, and that the Department's affirmative recruitment and 44\% Black recruit share negated any inference of discriminatory intent (\href{https://supreme.justia.com/cases/federal/us/426/229/}{Justia}). He warned that constitutionalizing disparate impact would threaten countless neutral tax, welfare, regulatory, and licensing laws, leaving any expansion of disparate-impact liability to legislative prescription (\href{https://supreme.justia.com/cases/federal/us/426/229/}{Justia}).

In Part III, Justice White reached the statutory claims under 42 U.S.C. \S~1981, D.C. Code \S~1-320, and 5 U.S.C. \S~3304(a)(1), holding that the District Court's finding of a positive correlation between Test 21 and police academy training scores---supported by the Futransky validation study, expert affidavits, and Civil Service Commission regulations---was sufficient to validate Test 21 without requiring a remand (\href{https://supreme.justia.com/cases/federal/us/426/229/}{Justia}).

\subsection{Concurring Opinion (Justice Stevens)}

Justice Stevens joined the Court's opinion (except its disapproval of the cases in footnote 12) but wrote separately to soften the sharp dichotomy between purpose and impact (\href{https://supreme.justia.com/cases/federal/us/426/229/}{Justia}). He observed that in collective governmental decisionmaking, ``the most probative evidence of intent will be objective evidence of what actually happened, rather than evidence describing the subjective state of mind of the actor,'' because actors are normally presumed to intend the natural consequences of their deeds (\href{https://supreme.justia.com/cases/federal/us/426/229/}{Justia}). Conversely, an otherwise legitimate law should not be struck down merely because an improper motive influenced one participant's vote (``A law conscripting clerics should not be invalidated because an atheist voted for it'') (\href{https://supreme.justia.com/cases/federal/us/426/229/}{Justia}).

On the merits, Justice Stevens declined to rely on the Department's good-faith recruitment efforts; instead, he found Test 21 constitutional because (1) requiring a uniform minimum level of literacy is manifestly relevant to policing, and (2) Test 21 was administered nationwide across the federal civil service, making the D.C. Police Department applicant pool too small a sample to overcome the presumption of neutrality (\href{https://supreme.justia.com/cases/federal/us/426/229/}{Justia}). On Part III, he agreed that training-program validation was legally permissible under 42 U.S.C. \S~1981 and D.C. Code \S~1-320 and not clearly erroneous under Fed. R. Civ. P. 52(a) (\href{https://supreme.justia.com/cases/federal/us/426/229/}{Justia}).

\subsection{Dissenting Opinion (Justice Brennan, joined by Justice Marshall)}

Justice Brennan did not address the constitutional holding in Part II, focusing instead on Part III (\href{https://supreme.justia.com/cases/federal/us/426/229/}{Justia}). First, he argued that the Court should not have decided the statutory question at all, because respondents' summary judgment motion rested purely on the Fifth Amendment and their complaint asserted no claim under 5 U.S.C. \S~3304 (\href{https://supreme.justia.com/cases/federal/us/426/229/}{Justia}). He also criticized footnote 12 for effectively pre-judging \textit{Arlington Heights}, which was already set for plenary argument (\href{https://supreme.justia.com/cases/federal/us/426/229/}{Justia}).

Reaching the statutory merits anyway, Justice Brennan argued that petitioners failed to meet the job-relatedness standards of \textit{Griggs}, \textit{Albemarle Paper Co. v. Moody}, 422 U.S. 405 (1975), EEOC Guidelines (29 C.F.R. \S\S~1607.4(c), 1607.5(b)(3)), and even the Civil Service Commission's own Supreme Court brief (\href{https://supreme.justia.com/cases/federal/us/426/229/}{Justia}). The record did not even contain the training-school examinations, every recruit eventually passed the training course with remedial help, and the Futransky study found no positive correlation between Test 21 (or training averages) and actual job performance for Black officers (\href{https://supreme.justia.com/cases/federal/us/426/229/}{Justia}). Validating one written verbal test merely by correlating it with another written verbal test in a training course risked rewarding general verbal facility rather than ``job-specific ability'' (\href{https://supreme.justia.com/cases/federal/us/426/229/}{Justia}).

\section{Examples: Future Applications}

\begin{enumerate}
\item \textbf{Same-Side (No Equal Protection Violation --- State Licensing Fee):} A state doubles its annual occupational licensing fee for commercial barbers and cosmetologists to fund health-inspection programs. Statistical studies show that Black and Hispanic entrepreneurs are disproportionately burdened by the higher fee and exit the industry at twice the rate of white licensees. Under \textit{Davis}, because the fee is facially neutral, serves a legitimate regulatory and fiscal purpose, and was enacted without racially discriminatory intent, it receives rational basis review and survives an Equal Protection challenge (\href{https://supreme.justia.com/cases/federal/us/426/229/}{Justia}).

\item \textbf{Same-Side (No Equal Protection Violation --- Civil Service Preference):} A state statute grants an absolute lifetime civil service hiring preference to military veterans. Because historical military demographics mean over 98\% of beneficiaries are men, the statute severely disadvantages female applicants for state jobs. Following \textit{Davis} (and later \textit{Personnel Administrator of Massachusetts v. Feeney}, 442 U.S. 256 (1979)), foreseeable disproportionate impact alone does not prove discriminatory purpose; because the legislature adopted the preference to reward veterans rather than ``because of'' its adverse effect on women, the statute survives constitutional scrutiny (\href{https://www.columbialawreview.org/content/process-scrutiny-motivational-inquiry-and-constitutional-rights/}{Columbia Law Review}; \href{https://scholarship.law.bu.edu/cgi/viewcontent.cgi?article=1823&context=faculty_scholarship}{BU Law}).

\item \textbf{Opposite-Side (Equal Protection Violation --- Covertly Discriminatory Municipal Rezoning):} A predominantly white suburb abruptly rezones a parcel from multifamily to single-family housing immediately after a nonprofit developer proposes racially integrated affordable apartments, departing sharply from its standard zoning procedures while council members make racially coded remarks in public hearings. Under \textit{Davis}'s ``totality of the relevant facts'' principle (elaborated in \textit{Village of Arlington Heights v. Metropolitan Housing Development Corp.}, 429 U.S. 252 (1977)), the procedural departures, sequence of events, and decisionmaker statements establish discriminatory purpose, triggering strict scrutiny (\href{https://supreme.justia.com/cases/federal/us/426/229/}{Justia}; \href{https://harvardlawreview.org/print/vol-139/presidential-speech-and-the-discrimination-paradox/}{Harvard Law Review}).

\item \textbf{Opposite-Side (Equal Protection Violation --- Selective Administrative Enforcement):} A city enforces a facially neutral ordinance requiring permits to operate street-food carts by granting permits to 95\% of white applicants while denying 98\% of otherwise identically situated Asian American applicants without any nonracial explanation. As Justice White recognized in citing \textit{Yick Wo v. Hopkins}, 118 U.S. 356 (1886), extreme, unexplained unequal administration of a facially neutral law permits an inference of intentional racial discrimination and violates the Equal Protection Clause (\href{https://supreme.justia.com/cases/federal/us/426/229/}{Justia}).

\item \textbf{Fence-Sitter (Ambiguous Intent --- Algorithmic Risk Scoring in Pretrial Detention):} A county adopts a facially neutral machine-learning risk-assessment algorithm to guide bail recommendations. The algorithm excludes race as an explicit variable but relies heavily on ZIP codes and prior arrest histories that strongly correlate with race, yielding significantly higher detention rates for Black defendants; county officials were warned of the racial disparity before adoption but proceeded anyway to reduce administrative costs. Under \textit{Davis}, mere awareness of a foreseeable racial disparity does not establish discriminatory purpose, pushing the case toward upholding the policy; yet under Justice Stevens's objective-evidence approach and \textit{Davis}'s totality-of-the-facts standard, plaintiffs could argue that choosing proxy variables known to mirror race despite readily available, equally accurate alternatives crosses the line into purposeful discrimination (\href{https://supreme.justia.com/cases/federal/us/426/229/}{Justia}; \href{https://www.columbialawreview.org/content/process-scrutiny-motivational-inquiry-and-constitutional-rights/}{Columbia Law Review}).
\end{enumerate}

\section{Critique}

\textit{Washington v. Davis} is one of the most consequential and contested decisions in modern equal protection jurisprudence, marking a decisive pivot from an \textit{antisubordination} (or group-disadvantaging) vision of equality toward an \textit{anticlassification} and fault-based intent model (\href{https://law.yale.edu/sites/default/files/documents/pdf/Faculty/Siegel_TheSupremeCourt2012TermForwardEqualityDivided.pdf}{Reva Siegel, Harvard Law Review}).

\textbf{Progressive and Doctrinal Critiques:} Progressive scholars and critical race theorists have mounted sustained attacks on \textit{Davis}'s intent requirement. In his landmark 1987 article, Charles R. Lawrence III argued that requiring proof of conscious discriminatory motivation ``ignores much of what we understand about how the human mind works,'' because much contemporary racial bias operates through unconscious racism and culturally shared stereotypes rather than overt animus (\href{https://scholarship.law.gwu.edu/cgi/viewcontent.cgi?params=/context/faculty_publications/article/2030/&path_info=Was_the_Disparate_Impact_Theory_a_Mistake.pdf}{GW Law}; \href{https://scholarship.law.georgetown.edu/cgi/viewcontent.cgi?article=1340&context=facpub}{Georgetown Law}). Similarly, Alan David Freeman criticized \textit{Davis} for enshrining a ``perpetrator perspective'' that treats racial inequality as a series of isolated, malicious acts rather than a systemic condition experienced by victims (\href{https://scholarship.law.gwu.edu/cgi/viewcontent.cgi?params=/context/faculty_publications/article/2030/&path_info=Was_the_Disparate_Impact_Theory_a_Mistake.pdf}{GW Law}). Owen Fiss contended that \textit{Davis} artificially downgraded \textit{Griggs} into a purely statutory rule despite its constitutional roots, while Reva Siegel documented how \textit{Davis} and its progeny created ``equality divided''---applying deferential review to minority challenges against facially neutral policies while applying strict scrutiny to race-conscious remedies (\href{https://digitalcommons.law.mercer.edu/cgi/viewcontent.cgi?article=1314&context=jour_mlr}{Mercer Law Review}; \href{https://law.yale.edu/sites/default/files/documents/pdf/Faculty/Siegel_TheSupremeCourt2012TermForwardEqualityDivided.pdf}{Reva Siegel, Harvard Law Review}). Critics also note a perverse incentive: \textit{Davis} encourages savvy policymakers to cloak discriminatory motives in facially neutral language, making constitutional violations exceedingly difficult to prove (\href{https://www.columbialawreview.org/content/process-scrutiny-motivational-inquiry-and-constitutional-rights/}{Columbia Law Review}).

\textbf{Originalist, Textualist, and Institutionalist Perspectives:} Defenders of \textit{Davis}---including institutionalist scholars like Paul Brest and many originalists and textualists---argue that the Equal Protection Clause prohibits government from purposefully denying equal legal protection, not from enacting neutral laws that interact unevenly with preexisting socioeconomic disparities (\href{https://scholarship.law.gwu.edu/cgi/viewcontent.cgi?params=/context/faculty_publications/article/2030/&path_info=Was_the_Disparate_Impact_Theory_a_Mistake.pdf}{GW Law}; \href{https://www.law.georgetown.edu/georgetown-law-journal/wp-content/uploads/sites/26/2021/11/Bernick_Antisubjugation.pdf}{Georgetown Law Journal}). Originalists such as Alfred Avins, Earl Maltz, John Harrison, and Christopher Green have emphasized that the original public meaning of ``equal protection of the laws'' focused on equal legal status and protection against discriminatory state laws (such as the Black Codes) rather than guaranteeing proportional group outcomes (\href{https://www.law.georgetown.edu/georgetown-law-journal/wp-content/uploads/sites/26/2021/11/Bernick_Antisubjugation.pdf}{Georgetown Law Journal}; \href{https://scholarship.law.duke.edu/cgi/viewcontent.cgi?article=1069&context=dlj_online}{Duke Law Journal Online}). From a separation-of-powers and democratic-accountability standpoint, Justice White's institutional concern carries weight: subjecting every neutral tax, fee, or competency exam with a racial disparity to strict scrutiny would transfer vast policymaking authority from elected legislatures to federal judges, whereas Congress remains free under Section 5 of the Fourteenth Amendment or Title VII to enact targeted disparate-impact remedies (\href{https://supreme.justia.com/cases/federal/us/426/229/}{Justia}).

\textbf{Underlying Competing Values:} At its core, \textit{Davis} pits \textit{formal equality, individual culpability, and judicial restraint} against \textit{substantive equality, dismantling structural racial stratification, and remedying unconscious or entrenched institutional bias}.

\section{Key Quotations}

\begin{quote}
``The central purpose of the Equal Protection Clause of the Fourteenth Amendment is the prevention of official conduct discriminating on the basis of race. It is also true that the Due Process Clause of the Fifth Amendment contains an equal protection component prohibiting the United States from invidiously discriminating between individuals or groups. \textit{Bolling v. Sharpe}, 347 U. S. 497 (1954). But our cases have not embraced the proposition that a law or other official act, without regard to whether it reflects a racially discriminatory purpose, is unconstitutional solely because it has a racially disproportionate impact.''
\end{quote}
\noindent--- Justice White, Opinion of the Court, 426 U.S. at 239 (\href{https://supreme.justia.com/cases/federal/us/426/229/}{Justia}; \href{https://www.courtlistener.com/opinion/109469/washington-v-davis/}{CourtListener}).

\begin{quote}
``Disproportionate impact is not irrelevant, but it is not the sole touchstone of an invidious racial discrimination forbidden by the Constitution. Standing alone, it does not trigger the rule, \textit{McLaughlin v. Florida}, 379 U. S. 184 (1964), that racial classifications are to be subjected to the strictest scrutiny, and are justifiable only by the weightiest of considerations.''
\end{quote}
\noindent--- Justice White, Opinion of the Court, 426 U.S. at 242 (\href{https://supreme.justia.com/cases/federal/us/426/229/}{Justia}; \href{https://www.courtlistener.com/opinion/109469/washington-v-davis/}{CourtListener}).

\begin{quote}
``A rule that a statute designed to serve neutral ends is nevertheless invalid, absent compelling justification, if in practice it benefits or burdens one race more than another would be far reaching and would raise serious questions about, and perhaps invalidate, a whole range of tax, welfare, public service, regulatory, and licensing statutes that may be more burdensome to the poor and to the average black than to the more affluent white.''
\end{quote}
\noindent--- Justice White, Opinion of the Court, 426 U.S. at 248 (\href{https://supreme.justia.com/cases/federal/us/426/229/}{Justia}; \href{https://www.courtlistener.com/opinion/109469/washington-v-davis/}{CourtListener}).

\begin{quote}
``Frequently the most probative evidence of intent will be objective evidence of what actually happened, rather than evidence describing the subjective state of mind of the actor. For, normally, the actor is presumed to have intended the natural consequences of his deeds.''
\end{quote}
\noindent--- Justice Stevens, Concurring Opinion, 426 U.S. at 253 (\href{https://supreme.justia.com/cases/federal/us/426/229/}{Justia}; \href{https://www.courtlistener.com/opinion/109469/washington-v-davis/}{CourtListener}).

\begin{quote}
``Where employers try to validate written qualification tests by proving a correlation with written examinations in a training course, there is a substantial danger that people who have good verbal skills will achieve high scores on both tests due to verbal ability, rather than `job-specific ability.'''
\end{quote}
\noindent--- Justice Brennan (joined by Justice Marshall), Dissenting Opinion, 426 U.S. at 270 (\href{https://supreme.justia.com/cases/federal/us/426/229/}{Justia}; \href{https://www.courtlistener.com/opinion/109469/washington-v-davis/}{CourtListener}).

\end{document}
